% W5i76_HardProblems.tex
% DFD 20-Agent Research Programme --- Iteration 76
% Hard Problems, Honest Negatives, and Open Questions

\documentclass[11pt,a4paper]{article}
\usepackage[utf8]{inputenc}
\usepackage{amsmath,amssymb}
\usepackage{booktabs}
\usepackage{geometry}
\usepackage{xcolor}
\usepackage{enumitem}
\geometry{margin=2.5cm}

\title{W5i76: Hard Problems and Honest Negatives\\[6pt]
\large N-Body Finite-Volume $\cdot$ $E_G$ Sensitivity Gaps $\cdot$ Prediction Horizon}
\author{DFD 20-Agent Research Programme}
\date{2026-03-24}

\begin{document}
\maketitle

\section{Hard Problem 1: N-Body Finite-Volume Effect}

\textbf{Status}: IDENTIFIED, NOT CRITICAL.

The $L = 500\,h^{-1}$Mpc simulation box shows $\sim 2\%$ sample variance at $k < 0.02\,h$/Mpc. The BAO feature at $k \sim 0.06\,h$/Mpc is resolved but with $\sim 1\%$ noise. Consequences:
\begin{itemize}[nosep]
\item DFD BAO-scale predictions cannot be validated at better than $\sim 1\%$ with current box size.
\item A $L = 1000\,h^{-1}$Mpc box ($\sim 800$k CPU-hours) is needed for BAO-scale DFD N-body validation.
\item \textbf{Not a RED risk}: the primary signal ($k = 0.1$--$0.5$) is robust and converged.
\end{itemize}

\textbf{Mitigation}: Include finite-volume caveat in N-body paper. Request larger allocation in Phase 2 HPC proposal.

\section{Hard Problem 2: $E_G$ Sensitivity Gap at $z > 0.7$}

At $z > 0.7$, DFD's $E_G$ excess over GR exceeds 4\%, but statistical errors grow rapidly due to reduced source galaxy density. Fisher forecast gives:
\begin{itemize}[nosep]
\item $z = 0.3$: $E_G^{\rm DFD} - E_G^{\rm GR} = 0.010$, $\sigma_{E_G} = 0.006$. Detection: $1.7\sigma$.
\item $z = 0.5$: $E_G^{\rm DFD} - E_G^{\rm GR} = 0.014$, $\sigma_{E_G} = 0.008$. Detection: $1.8\sigma$.
\item $z = 0.7$: $E_G^{\rm DFD} - E_G^{\rm GR} = 0.016$, $\sigma_{E_G} = 0.012$. Detection: $1.3\sigma$.
\item Combined $z < 0.7$: $3.2\sigma$ (Fisher optimal).
\end{itemize}

The high-$z$ sensitivity gap means $E_G$ alone cannot distinguish DFD from GR at $z > 0.7$. Combined probes (lensing + clustering + ISW) are needed.

\textbf{Honest negative}: $E_G$ is a $3\sigma$ test, not a $5\sigma$ discovery channel. DFD's strongest near-term falsification comes from GW lensing (pass/fail), not $E_G$ (statistical).

\section{Hard Problem 3: Prediction Horizon and Theory-Data Gap}

Of the 28 compiled predictions:
\begin{itemize}[nosep]
\item 5 testable within 2 years (imminent)
\item 8 testable within 5 years (BICEP Array, LiteBIRD, JUNO, CMB-S4 Phase 1)
\item 10 testable within 10 years (Einstein Telescope, LISA, next-gen surveys)
\item 5 testable only with future technology ($\dot{G}/G$ at $10^{-15}$/yr, $\alpha$ running at $10^{-8}$)
\end{itemize}

The 5 long-horizon predictions are \textit{structurally unfalsifiable on decade timescales}. This is an honest limitation, not a flaw: most unified theories have similar or worse horizons. DFD's 5 imminent tests are exceptional by comparison.

\section{Hard Problem 4: $\alpha$ Residual (Persistent YELLOW)}

The $\alpha$ residual stands at 0.55 ppm (0.006 ppm from CODATA, but theoretical uncertainty from Toeplitz truncation is 0.55 ppm). Reducing this requires:
\begin{itemize}[nosep]
\item Computing the $n = 5$ spectral action coefficient (currently only $n \leq 4$ computed)
\item Modular form verification to confirm $k = 60$ (currently verified to $k = 58$ by direct computation)
\item Neither is feasible without dedicated symbolic computation ($\sim 10^6$ CPU-hours for $n = 5$)
\end{itemize}

\textbf{Classification}: Permanent YELLOW. Not a threat to the programme (sub-ppm accuracy is already unprecedented), but prevents promotion to GREEN.

\section{Hard Problem 5: Neutrino Mass Hierarchy Assumption}

DFD predicts $\Sigma m_\nu = 61.4$ meV assuming normal ordering. If inverted ordering is established by JUNO, the DFD prediction shifts to $\Sigma m_\nu = 101$ meV, which is:
\begin{itemize}[nosep]
\item Consistent with current cosmological bounds ($< 120$ meV at 95\% CL)
\item But removes the distinctive ``minimal normal ordering'' prediction
\item DFD's geometry (spin$^c$ bundle on $\CP^2$) slightly favours normal ordering (3 light eigenvalues $\neq$ 2 light)
\end{itemize}

\textbf{Not a falsification risk}: DFD accommodates both orderings. But inverted ordering would weaken the predictive claim. Honest negative logged.

\section{Open Questions from i76}

\begin{enumerate}[nosep]
\item Can the $L = 1000\,h^{-1}$Mpc N-body box be run within the current HPC allocation? (Requires $\sim 800$k additional CPU-hours.)
\item What is the optimal $E_G$ binning strategy for maximising DFD--GR discrimination?
\item Can the multi-messenger paper include specific ET sensitivity curves (currently only O5)?
\item Is there a way to test the DFD halo mass function prediction with current eROSITA data?
\item What is the systematic floor for $D_L^{\rm GW}/D_L^{\rm EM}$ measurements in O5?
\end{enumerate}

\section{Summary}

5 hard problems identified. 0 are RED risks. 2 are permanent YELLOWs (already known). 1 is a finite-volume technical issue (solvable). 2 are honest limitations logged for intellectual completeness. Programme integrity: \textbf{MAINTAINED}.

\end{document}
